\chapter{Lists, Stacks, Queues}\label{chap:linear}

This chapter describes the implementation of lists, stacks, and queues, which are derived from the fundamental deque and array data structures.

\section{List interface}\label{lists:interface}

Our list implementation is a direct wrapper around the deque implementation. Since a deque supports all list operations (insertion/removal at both ends, iteration), we simply expose the deque functionality with list-specific terminology. This avoids code duplication and ensures consistency.

\section{Stack interface}\label{stacks:interface}

Stacks (LIFO - Last In, First Out) are implemented by specializing either dynamic arrays or deques. The user can choose the underlying implementation based on their specific needs (e.g., cache locality of arrays vs. pointer stability of deques).

The stack interface provides standard operations: \texttt{push}, \texttt{pop}, \texttt{top}, and \texttt{is\_empty}.

\section{Queue interface}\label{queues:interface}

Queues (FIFO - First In, First Out) are also implemented by specializing deques. While dynamic arrays could be used, removing elements from the front of an array is inefficient ($O(N)$) unless a circular buffer approach is used. Deques naturally support $O(1)$ insertion at the back and removal from the front, making them an ideal choice for queues.

The queue interface provides standard operations: \texttt{enqueue} (push back), \texttt{dequeue} (pop front), \texttt{front}, and \texttt{is\_empty}.

\section{Benchmark}

We benchmarked the performance of stacks implemented with dynamic arrays versus deques.

\subsection{Stacks from dynamic arrays \textit{vs.} from deques}

Dynamic arrays generally offer better cache locality and less memory overhead per element compared to deques (which require two pointers per element). However, dynamic arrays incur occasional resizing costs. Deques provide consistent $O(1)$ performance for operations but have higher memory overhead and poorer cache locality due to non-contiguous memory allocation.

\subsection{Observations and conclusions}

For most use cases where the maximum size is not known in advance and performance is critical, dynamic arrays are often preferred due to cache efficiency. However, if pointer stability is required (i.e., addresses of elements must not change), deques are the better choice.
